\documentclass[a4paper,12pt]{article}

\usepackage[utf8]{inputenc}
\usepackage[T1]{fontenc}
\usepackage[ngerman]{babel}
\usepackage{geometry}
\usepackage{titlesec}
\usepackage{enumitem}
\usepackage{booktabs}
\usepackage{array}
\usepackage{tabularx}
\usepackage{parskip}
\usepackage{hyperref}
\usepackage{lmodern}
\usepackage{microtype}
\usepackage{amssymb}
\usepackage[simplified]{pgf
-umlcd}

\geometry{
  a4paper,
  left=2.5cm,
  right=2.5cm,
  top=3cm,
  bottom=3cm
}

% Section formatting
\titleformat{\section}[block]{\large\bfseries}{Kapitel \thesection:}{0.5em}{}
\titleformat{\subsection}[block]{\normalsize\bfseries}{\thesubsection}{0.5em}{}
\titleformat{\subsubsection}[block]{\normalsize\bfseries}{\thesubsubsection}{0.5em}{}

\setcounter{secnumdepth}{3}
\setcounter{tocdepth}{3}

% Custom colors / highlighting for placeholder text
\usepackage{xcolor}
\newcommand{\placeholder}[1]{\textcolor{gray}{[#1]}}


% List settings
\renewcommand{\labelitemiii}{$\blacksquare$}

\setlist[itemize,1]{label=\textbullet, leftmargin=1.5em}
\setlist[itemize,2]{label=\textopenbullet, leftmargin=1.5em}
\setlist[itemize,3]{label=$\blacksquare$, leftmargin=1.5em}
\setlist[itemize,4]{label=\textbullet, leftmargin=1.5em}

\begin{document}

% ============================================================
% Title Page
% ============================================================
\begin{titlepage}
    \centering
    \vspace*{4cm}

    {\LARGE\bfseries Programmentwurf \placeholder{Bezeichnung}\par}

    \vspace{3cm}

    \begin{tabular}{ll}
        \textbf{Name:}           & \placeholder{Name, Vorname} \\[0.5em]
        \textbf{Matrikelnummer:} & \placeholder{MNR}           \\[0.5em]
        \textbf{Abgabedatum:}    & \placeholder{DATUM}
    \end{tabular}
\end{titlepage}

% ============================================================
% Allgemeine Anmerkungen
% ============================================================
\newpage
\subsection*{Allgemeine Anmerkungen:}

\begin{itemize}
    \item es darf nicht auf andere Kapitel als Leistungsnachweis verwiesen werden
          (z.\,B. in der Form \glqq{}XY wurde schon in Kapitel 2 behandelt, daher hier keine
          Ausführung\grqq{})
    \item alles muss in UTF-8 codiert sein (Text und Code)
    \item sollten mündliche Aussagen den schriftlichen Aufgaben widersprechen, gelten die
          schriftlichen Aufgaben (ggf.\ an Anpassung der schriftlichen Aufgaben erinnern!)
    \item alles muss ins Repository (Code, Ausarbeitung und alles was damit zusammenhängt)
    \item die Beispiele sollten wenn möglich vom aktuellen Stand genommen werden
          \begin{itemize}
              \item finden sich dort keine entsprechenden Beispiele, dürfen auch ältere Commits
                    unter Verweis auf den Commit verwendet werden
              \item Ausnahme: beim Kapitel \glqq{}Refactoring\grqq{} darf von vorne herein aus
                    allen Ständen frei gewählt werden (mit Verweis auf den entsprechenden Commit)
          \end{itemize}
    \item falls verlangte Negativ-Beispiele nicht vorhanden sind, müssen entsprechend mehr
          Positiv-Beispiele gebracht werden
          \begin{itemize}
              \item Achtung: werden im Code entsprechende Negativ-Beispiele gefunden, gibt es
                    keine Punkte für die zusätzlichen Positiv-Beispiele
              \item Beispiele
                    \begin{itemize}
                        \item \glqq{}Nennen Sie jeweils eine Klasse, die das SRP einhält bzw.\ verletzt.\grqq{}
                              \begin{itemize}
                                  \item Antwort: Es gibt keine Klasse, die SRP verletzt, daher hier 2 Klassen,
                                        die SRP einhalten: [Klasse 1], [Klasse 2]
                                  \item Bewertung: falls im Code tatsächlich keine Klasse das SRP verletzt:
                                        volle Punktzahl ODER falls im Code mind.\ eine Klasse SRP verletzt:
                                        halbe Punktzahl
                              \end{itemize}
                    \end{itemize}
          \end{itemize}
    \item verlangte Positiv-Beispiele müssen gebracht werden
    \item Code-Beispiel = Code in das Dokument kopieren
\end{itemize}

\section{Testen von TeX}
\subsection{Link zur Dokumentation vom UML-Package}
https://ctan.uib.no/graphics/pgf/contrib/pgf-umlcd/pgf-umlcd-manual.pdf
\subsection{Beispiel für UML-Diagramm}
\begin {tikzpicture}
\begin{class}[ text width =7 cm ]{ Class }{0 ,0}
    \attribute{+ Public }
    \attribute{\# Pr o te ct e d }
    \attribute{ - Private }
    \attribute{$\ sim $ Package }
\end{class}
\begin{class}[ text width =7 cm ]{ B a n k A c c o u n t }{0 , -3}
    \attribute{+ owner : String }
    \attribute{+ balance : Dollars }
    \operation{+ deposit ( amount : Dollars ) }
    \operation{+ w i t h d r a w a l ( amount : Dollars ) }
    \operation {\# u p d a t e B a l a n c e ( n e w B a l a n c e : Dollars
        ) }
\end{class}
\end{tikzpicture}


% ============================================================
% Kapitel 1: Einführung
% ============================================================
\section{Einführung}

\subsection*{Übersicht über die Applikation}

\placeholder{Was macht die Applikation? Wie funktioniert sie? Welches Problem löst sie/welchen Zweck hat sie?}


\subsection*{Wie startet man die Applikation?}

\placeholder{Wie startet man die Applikation? Welche Voraussetzungen werden benötigt? Schritt-für-Schritt-Anleitung}

\subsection*{Wie testet man die Applikation?}

\placeholder{Wie testet man die Applikation? Welche Voraussetzungen werden benötigt? Schritt-für-Schritt-Anleitung}

% ============================================================
% Kapitel 2: Clean Architecture
% ============================================================
\section{Clean Architecture}

\subsection*{Was ist Clean Architecture?}

\placeholder{allgemeine Beschreibung der Clean Architecture in eigenen Worten}

\subsection*{Analyse der Dependency Rule}

\placeholder{(1 Klasse, die die Dependency Rule einhält und eine Klasse, die die Dependency Rule verletzt);
    jeweils UML der Klasse und Analyse der Abhängigkeiten in beide Richtungen (d.\,h., von wem hängt die Klasse
    ab und wer hängt von der Klasse ab) in Bezug auf die Dependency Rule}

\subsubsection*{Positiv-Beispiel: Dependency Rule}

\subsubsection*{Negativ-Beispiel: Dependency Rule}

\subsection*{Analyse der Schichten}

\placeholder{jeweils 1 Klasse zu 2 unterschiedlichen Schichten der Clean-Architecture: jeweils UML der Klasse
    (ggf.\ auch zusammenspielenden Klassen), Beschreibung der Aufgabe, Einordnung mit Begründung in die Clean-Architecture}

\subsubsection*{Schicht: \placeholder{Name}}

\subsubsection*{Schicht: \placeholder{Name}}

% ============================================================
% Kapitel 3: SOLID
% ============================================================
\section{SOLID}

\subsection*{Analyse Single-Responsibility-Principle (SRP)}

\placeholder{jeweils eine Klasse als positives und negatives Beispiel für SRP;
    jeweils UML der Klasse und Beschreibung der Aufgabe bzw.\ der Aufgaben und möglicher Lösungsweg
    des Negativ-Beispiels (inkl.\ UML)}

\subsubsection*{Positiv-Beispiel}

\subsubsection*{Negativ-Beispiel}

\newpage

\subsection*{Analyse Open-Closed-Principle (OCP)}

\placeholder{jeweils eine Klasse als positives und negatives Beispiel für OCP;
    jeweils UML der Klasse und Analyse mit Begründung, warum das OCP erfüllt/nicht erfüllt wurde –
    falls erfüllt: warum hier sinnvoll/welches Problem gab es? Falls nicht erfüllt: wie könnte man
    es lösen (inkl.\ UML)?}

\subsubsection*{Positiv-Beispiel}

\subsubsection*{Negativ-Beispiel}

\bigskip

\subsection*{Analyse Liskov-Substitution- (LSP), Interface-Segregation- (ISP),\\
    Dependency-Inversion-Principle (DIP)}

\placeholder{jeweils eine Klasse als positives und negatives Beispiel für entweder LSP oder ISP oder DIP);
    jeweils UML der Klasse und Begründung, warum man hier das Prinzip erfüllt/nicht erfüllt wird}

\smallskip
\noindent\placeholder{Anm.: es darf nur ein Prinzip ausgewählt werden; es darf NICHT z.\,B.\ ein positives
    Beispiel für LSP und ein negatives Beispiel für ISP genommen werden}

\subsubsection*{Positiv-Beispiel}

\subsubsection*{Negativ-Beispiel}

% ============================================================
% Kapitel 4: Weitere Prinzipien
% ============================================================
\section{Weitere Prinzipien}

\subsection*{Analyse GRASP: Geringe Kopplung}

\placeholder{jeweils eine bis jetzt noch nicht behandelte Klasse als positives und negatives Beispiel
    geringer Kopplung; jeweils UML Diagramm mit zusammenspielenden Klassen, Aufgabenbeschreibung und
    Begründung für die Umsetzung der geringen Kopplung bzw.\ Beschreibung, wie die Kopplung aufgelöst
    werden kann}

\subsubsection*{Positiv-Beispiel}

\subsubsection*{Negativ-Beispiel}

\bigskip

\subsection*{Analyse GRASP: Hohe Kohäsion}

\placeholder{eine Klasse als positives Beispiel hoher Kohäsion; UML Diagramm und Begründung, warum
    die Kohäsion hoch ist}

\subsection*{Don't Repeat Yourself (DRY)}

\placeholder{ein Commit angeben, bei dem duplizierter Code/duplizierte Logik aufgelöst wurde;
    Code-Beispiele (vorher/nachher); begründen und Auswirkung beschreiben}

% ============================================================
% Kapitel 5: Unit Tests
% ============================================================
\newpage
\section{Unit Tests}

\subsection*{10 Unit Tests}

\placeholder{Nennung von 10 Unit-Tests und Beschreibung, was getestet wird}

\bigskip
\begin{tabularx}{\textwidth}{|l|X|l|}
    \hline
    \textbf{Unit Test} & \textbf{Beschreibung} & \textbf{Klasse\#Methode} \\
    \hline
                       &                       &                          \\[1.2em]
    \hline
                       &                       &                          \\[1.2em]
    \hline
                       &                       &                          \\[1.2em]
    \hline
                       &                       &                          \\[1.2em]
    \hline
                       &                       &                          \\[1.2em]
    \hline
                       &                       &                          \\[1.2em]
    \hline
                       &                       &                          \\[1.2em]
    \hline
                       &                       &                          \\[1.2em]
    \hline
                       &                       &                          \\[1.2em]
    \hline
                       &                       &                          \\[1.2em]
    \hline
\end{tabularx}

\bigskip

\subsection*{ATRIP: Automatic}

\placeholder{Begründung/Erläuterung, wie 'Automatic' realisiert wurde}

\subsection*{ATRIP: Thorough}

\placeholder{jeweils 1 positives und negatives Beispiel zu 'Thorough'; jeweils Code-Beispiel,
    Analyse und Begründung, was professionell/nicht professionell ist}

\subsection*{ATRIP: Professional}

\placeholder{jeweils 1 positives und negatives Beispiel zu 'Professional'; jeweils Code-Beispiel,
    Analyse und Begründung, was professionell/nicht professionell ist}

\subsection*{Code Coverage}

\placeholder{Code Coverage im Projekt analysieren und begründen}

\subsection*{Fakes und Mocks}

\placeholder{Analyse und Begründung des Einsatzes von 2 Fake/Mock-Objekten;
    zusätzlich jeweils UML Diagramm der Klasse}

% ============================================================
% Kapitel 6: Domain Driven Design
% ============================================================
\newpage
\section{Domain Driven Design}

\subsection*{Ubiquitous Language}

\placeholder{4 Beispiele für die Ubiquitous Language; jeweils Bezeichnung, Bedeutung und kurze
    Begründung, warum es zur Ubiquitous Language gehört}

\bigskip
\begin{tabularx}{\textwidth}{|l|X|X|}
    \hline
    \textbf{Bezeichnung} & \textbf{Bedeutung} & \textbf{Begründung} \\
    \hline
                         &                    &                     \\[1.2em]
    \hline
                         &                    &                     \\[1.2em]
    \hline
                         &                    &                     \\[1.2em]
    \hline
                         &                    &                     \\[1.2em]
    \hline
\end{tabularx}

\bigskip

\subsection*{Entities}

\placeholder{UML, Beschreibung und Begründung des Einsatzes einer Entity; falls keine Entity
    vorhanden: ausführliche Begründung, warum es keines geben kann/hier nicht sinnvoll ist}

\subsection*{Value Objects}

\placeholder{UML, Beschreibung und Begründung des Einsatzes eines Value Objects; falls kein
    Value Object vorhanden: ausführliche Begründung, warum es keines geben kann/hier nicht sinnvoll ist}

\subsection*{Repositories}

\placeholder{UML, Beschreibung und Begründung des Einsatzes eines Repositories; falls kein
    Repository vorhanden: ausführliche Begründung, warum es keines geben kann/hier nicht sinnvoll ist}

\subsection*{Aggregates}

\placeholder{UML, Beschreibung und Begründung des Einsatzes eines Aggregates; falls kein
    Aggregate vorhanden: ausführliche Begründung, warum es keines geben kann/hier nicht sinnvoll ist}

% ============================================================
% Kapitel 7: Refactoring
% ============================================================
\section{Refactoring}

\subsection*{Code Smells}

\placeholder{jeweils 1 Code-Beispiel zu 2 Code Smells aus der Vorlesung; jeweils Code-Beispiel
    und einen möglichen Lösungsweg bzw.\ den genommenen Lösungsweg beschreiben (inkl.\ (Pseudo-)Code)}

\subsection*{2 Refactorings}

\placeholder{2 unterschiedliche Refactorings aus der Vorlesung anwenden, begründen, sowie UML
    vorher/nachher liefern; jeweils auf die Commits verweisen}

% ============================================================
% Kapitel 8: Entwurfsmuster
% ============================================================
\section{Entwurfsmuster}

\placeholder{2 unterschiedliche Entwurfsmuster aus der Vorlesung (oder nach Absprache auch andere)
    jeweils sinnvoll einsetzen, begründen und UML-Diagramm}

\subsection*{Entwurfsmuster: \placeholder{Name}}

\subsection*{Entwurfsmuster: \placeholder{Name}}

\end{document}